\section{Reciprocal symmetry of the real-emission integral}
After the first energy integration, write $s=1+H$, with $H>0$,
and $z=s+1-u-s/u$. The two algebraic boundaries of the first
integral are $x_-=s-u$ and $x_+=s-s/u$. Under $u\mapsto s/u$,
$z$ is unchanged and these boundaries exchange places. The
oriented boundary difference therefore changes sign. Its
Jacobian $j=1-s/u^2$ also changes sign, so the complete last
integrand obeys
\begin{equation}
 \frac{s}{u^2}f(s/u)=f(u).
\end{equation}
This identity holds separately for the regular rational,
energy-logarithm, $W$, and Bose-restoring sectors. It can also be
checked coefficient by coefficient: $L_1,L_2,L_3$ change sign,
whereas $L_1L_2$ is invariant under the reciprocal transformation.

The physical upper branch has $\sqrt{s}<u<s$. Its reflected
partner is $1<u<\sqrt{s}$. A common cut $z>\epsilon$ selects
reciprocal endpoints $u_-(\epsilon)u_+(\epsilon)=s$. Consequently
\begin{equation}
 \int_{\sqrt{s}}^{u_+(\epsilon)}f(u)\,\dd u
 =\frac12\int_{u_-(\epsilon)}^{u_+(\epsilon)}f(u)\,\dd u.
\end{equation}
Thus the square-root lower boundary can be eliminated exactly,
without changing the infrared subtraction or making an asymptotic
approximation.

Let $A$ be the coefficient of $1/(s-u)$ at the upper endpoint.
Reciprocal symmetry makes it also the coefficient of $1/(u-1)$
at the lower endpoint. Define
\begin{equation}
 \operatorname{FP}_{[1,s]}\int f
 =\lim_{\delta\to0^+}
 \left[\int_{1+\delta}^{s-\delta}f(u)\,\dd u+2A\ln\delta\right].
\end{equation}
Since $z=H(u-1)+O((u-1)^2)$ at the lower boundary and
$z=H(s-u)/s+O((s-u)^2)$ at the upper boundary, the finite part
in the original $z$ convention is
\begin{equation}
 F_z=\frac12\operatorname{FP}_{[1,s]}\int f
                  +A\ln\frac{H}{\sqrt{s}}.
\end{equation}
The explicit logarithm is essential. The replacement of the lower
boundary does not permit changing the finite subtraction constant.

The Bose-restoring term also has a double endpoint pole. If
\begin{equation}
 f(s-b)=D/b^2+A/b+O(1),
\end{equation}
reciprocal symmetry gives $f(1+a)=D/(sa^2)+A/a+O(1)$.
The finite-part definition then also subtracts
$D(1+1/s)/\delta$. Expanding both cut boundaries to second order
gives the corrected general formula
\begin{equation}
 F_z=\frac12\operatorname{FP}_{[1,s]}\int f
       +A\ln\frac{H}{\sqrt{s}}-\frac{D(s+1)}{2sH}.
\end{equation}
The previous expression is the case $D=0$. The extra constant
comes from the double pole multiplied by the second-order change
of the endpoint, and must be kept for the Bose term.

\paragraph{Eliminating the apparent threshold at $s=9$.}
The regular rational primitive contains ten words whose first letters are
\begin{equation}
 q_\pm=\frac{s-3\pm\sqrt{(s-1)(s-9)}}2.
\end{equation}
They can be replaced by a single real dilogarithm and GPLs whose letters
already occur in the other terms. This replacement follows from the
coefficients of our primitive. It does not use a total-cross-section formula.
Throughout this paragraph the GPL basepoint is $r=\sqrt{s}$, as in the
original primitive. Define
\begin{align}
 A(u)&=\frac{(u-1)(u+s)^2}{u(u+1)^2(s-u)},
 &Q(u)&=u^2-(s-3)u+s,\\
 (k_0,k_1,k_s,k_{-s},k_{-1})&=(-1,1,-1,2,-2),
 &C(s)&=\frac{128s(s+5)}{3(s-1)}.
\end{align}
The ten coefficients are precisely $C(s)k_b$ for each of the two first
letters $q_\pm$ and each $b\in\{0,1,s,-s,-1\}$. Since $A(r)=1$,
\begin{equation}
 L_3(u)=\ln A(u)=\sum_b k_bG(b;u).
\end{equation}
Factorization of $1-A$ gives the rational differential identity
\begin{equation}
 \frac{Q'(u)}{Q(u)}
 =\frac{\dd}{\dd u}\ln(1-A(u))
 -\frac1{u-r}-\frac1{u+r}
 +\frac1u+\frac1{u-s}+\frac2{u+1}.
\end{equation}
Here the logarithmic derivative denotes the rational quotient
$-A'(u)/(1-A(u))$ wherever its separate terms have removable poles.
Multiplying by $L_3$ and using
$\dd\operatorname{Li}_2(1-A)=-\ln A\,\dd\ln(1-A)$ yields
\begin{align}
 \sum_{\epsilon=\pm}\sum_b k_bG(q_\epsilon,b;u)
 ={}&-\operatorname{Li}_2(1-A(u))\\
 &+\sum_b k_b\bigl[-G(r,b;u)-G(-r,b;u)
       +G(0,b;u)\bigr]\\
 &+\sum_b k_b\bigl[G(s,b;u)+2G(-1,b;u)\bigr].
 \label{eq:comp-real-qpair-reduction}
\end{align}
Both sides vanish at $u=r$, which fixes the integration constant.
Although $r$ is a first letter on the right, the summed second-letter
function $L_3$ vanishes there, so the corresponding lower endpoint is
regular. For $r<u<s$, $A(u)>0$. Thus $1-A(u)<1$, and the dilogarithm
in this expression lies on its real branch. At a zero of $Q$, $A=1$
and this dilogarithm is regular. The formula also holds by continuity
at $s=9$, without choosing a side of an artificial threshold.

Replacing the ten words by \eqref{eq:comp-real-qpair-reduction} removes
$q_\pm$ from the primitive. It reintroduces the letters $\pm r$, so it
is a trade between two representations rather than a simultaneous
elimination of every auxiliary root. The accompanying coefficient
record contains the full replaced primitive. Its ten input coefficients,
rational differential identity, and lower-endpoint constant are checked
exactly by \texttt{compton\_real\_equivalence\_qpair.py}.

\paragraph{Combining the rational and $W$ primitives.}
The two primitives should be added before selecting a function basis.
Their original lists contain 48 and 19 nonempty GPL words. Collecting
identical words and canceling their exact rational coefficients leaves
43 words. In particular, the six Gram words discussed above disappear.
The coefficients in this combined list can be reduced further by shuffle
identities. For weight two,
\begin{equation}
 G(a,b)+G(b,a)=G(a)G(b),\qquad G(a,a)=\frac12G(a)^2.
\end{equation}
Applying these identities only to pairs present in the combined list leaves
40 GPL words and three elementary logarithm products.

There is a further reduction in the weight-three part. Define
\begin{equation}
 K=\frac{32s(2s+1)}{(s-1)^2},\qquad g=G(0;u)=\ln(u/\sqrt{s}).
\end{equation}
The entire weight-three contribution of the combined primitive is
\begin{align}
 K\bigl[&G(0,0,s)+G(0,s,0)-G(0,0,1)-G(0,1,0)
          -2G(0,0,0)\bigr]\\
 &=K\bigl[g\{G(0,s)-G(0,1)\}
          -G(0,0,s)+G(0,0,1)-g^3/3\bigr].
\end{align}
This follows from
$G(0)G(0,b)=2G(0,0,b)+G(0,b,0)$ for $b=1,s$.
Consequently the combined representation requires 37 nonempty GPL words,
three logarithm products, two logarithm--weight-two products, and one
logarithm cube. These are representation counts, not a claim of functional
independence of all remaining GPLs. The rational term is included throughout.

By contrast, substituting \eqref{eq:comp-real-qpair-reduction} before the
same combination gives 50 GPL words. Weight-two shuffle reduction leaves
42 words, eight logarithm products, and the additional dilogarithm.
The root-pair identity is therefore useful for displaying regularity at
$s=9$, but the original root-pair representation gives the shorter combined
formula. The full coefficient dictionary for the 37-word form and exact
reconstruction checks are recorded by
\texttt{compton\_real\_equivalence\_combine.py}.

\paragraph{Classical trilogarithms for the weight-three term.}
For positive real $z$, define $\mathcal L_n(z)$ here as
$\operatorname{Re}\operatorname{Li}_n(z+\ii0)$, for $n=2,3$.
The two boundary values are complex conjugates, so their real parts agree.
The differential relations are
\begin{equation}
 \frac{\dd\mathcal L_2(z)}{\dd z}=-\frac{\ln|1-z|}{z},
 \qquad
 \frac{\dd\mathcal L_3(z)}{\dd z}=\frac{\mathcal L_2(z)}{z}.
\end{equation}
They determine the continuation needed on each of the intervals used below.
For $b=1,s$ and $r<u<s$, direct differentiation and the zero value at $u=r$
give
\begin{align}
 G(0,b;u)&=-\mathcal L_2(u/b)+\mathcal L_2(r/b)
                      -\ln|1-r/b|\,g,\\
 G(0,0,b;u)&=-\mathcal L_3(u/b)+\mathcal L_3(r/b)
       +\mathcal L_2(r/b)g-\tfrac12\ln|1-r/b|\,g^2.
\end{align}
No singularity at $u=b$ is crossed along either open integration interval.
The complete weight-three contribution therefore reduces to
\begin{align}
 \mathcal F^{(3)}_{R_0+W}(u)=K\bigl[&
 \mathcal L_3(u/s)-\mathcal L_3(u)
 -\mathcal L_3(1/r)+\mathcal L_3(r)\\
 &+g\{\mathcal L_2(u)-\mathcal L_2(u/s)\}
 +\tfrac12g^2\ln r-\tfrac13g^3\bigr].
 \label{eq:comp-real-classical-weight-three}
\end{align}
The logarithmic coefficient uses
$|1-r|/|1-r/s|=r$. Thus the five original weight-three words require
only classical trilogarithms, dilogarithms, and logarithms. The lower
endpoint still vanishes exactly.

The remaining weight-two GPLs admit a uniform classical representation.
For distinct letters $a,b$ write
\begin{equation}
 z_a(u;b)=\frac{a-u}{a-b},\qquad z_a(r;b)=\frac{a-r}{a-b}.
\end{equation}
Then
\begin{align}
 G(a,b;u)={}&-\operatorname{Li}_2(z_a(u;b))
              +\operatorname{Li}_2(z_a(r;b))\\
 &-\operatorname{Log}(1-z_a(r;b))G(a;u).
 \label{eq:comp-real-classical-weight-two}
\end{align}
The logarithm and dilogarithm on the right must be continued together
from $u=r$ along the same path as the original GPL. More explicitly, their
logarithms satisfy
\begin{equation}
 \operatorname{Log}(1-z_a(u;b))
 -\operatorname{Log}(1-z_a(r;b))=G(b;u).
\end{equation}
This condition prevents independent principal-branch substitutions at the
two endpoints. Differentiation of \eqref{eq:comp-real-classical-weight-two}
then gives $G(b;u)/(u-a)$, and its initial value is zero. It therefore
proves the identity, including the continuation around removable letters.
Conjugate letters are combined before taking the real physical value.
Together with \eqref{eq:comp-real-classical-weight-three}, this eliminates
all GPL notation from the combined primitive if a classical-function
representation is preferred. Expanding every such identity separately
need not shorten the displayed expression, so conjugate pairs and common
coefficients should remain grouped.

At the upper endpoint the weight-three part simplifies further.
Put $\ell=\tfrac12\ln s$. The real inversion identities for $z>1$ are
\begin{align}
 \mathcal L_2(z)+\operatorname{Li}_2(z^{-1})
 &=\frac{\pi^2}{3}-\frac12\ln^2 z,\\
 \mathcal L_3(z)-\operatorname{Li}_3(z^{-1})
 &=\frac{\pi^2}{3}\ln z-\frac16\ln^3z.
\end{align}
Their derivatives agree on $z>1$, and their constants follow from
$\operatorname{Li}_2(1)=\pi^2/6$ and
$\operatorname{Li}_3(1)=\zeta(3)$. Applying them at $z=s,r$ gives
\begin{equation}
 \mathcal F^{(3)}_{R_0+W}(s)
 =K\left[\zeta(3)-\operatorname{Li}_3(s^{-1})
 -\ell\operatorname{Li}_2(s^{-1})
 -\frac{\pi^2\ell}{6}-\frac{2\ell^3}{3}\right].
 \label{eq:comp-real-classical-weight-three-endpoint}
\end{equation}
The lower endpoint is zero. The resulting expression depends only on
$s$ and contains no auxiliary square root or continued polylogarithm.

\paragraph{Finite energy endpoints before further combination.}
For the combined rational and $W$ primitive, cancellation of common
factors makes every rational coefficient regular at $u=s$. Thus its
finite upper endpoint needs only the coefficients evaluated at $s$.
No derivative of a coefficient multiplying a power-divergent GPL is
required. This statement follows by exact cancellation of their rational
denominators, whose endpoint factors are powers of $s$, $s-1$, and $s+3$.

For a one-letter word define its constant endpoint term by
\begin{equation}
 h_a=\begin{cases}
 \operatorname{Log}[(s-a)/(r-a)],&a\ne s,\\
 -\ln(s-r),&a=s.
 \end{cases}
\end{equation}
The second line removes the term $\ln(s-u)$ from $G(s;u)$.
For a weight-two word write $d_{ab}$ for the corresponding constant term.
For $a\ne s$, it is the boundary value of
\eqref{eq:comp-real-classical-weight-two}. For $a=s$ it is explicitly
\begin{equation}
 d_{sb}=\operatorname{Li}_2\left(\frac{s-r}{s-b}\right)
 +\operatorname{Log}\left(\frac{r-b}{s-b}\right)\ln(s-r).
\end{equation}
Indeed, the omitted divergent term is
$-\operatorname{Log}[(r-b)/(s-b)]\ln(s-u)$.
All logarithms use the continuation specified above. The only
logarithm products in the reduced primitive are $G(0)^2$, $G(0)G(1)$,
and $G(0)G(s)$, so their constant terms are simply the corresponding
products of the $h_a$. There is no product of two divergent one-letter
functions in this representation.

Replacing each word in the coefficient table by these constant endpoint
terms, and replacing its entire weight-three block by
\eqref{eq:comp-real-classical-weight-three-endpoint}, gives a function
$E_{R_0+W}(s)$ of energy alone. Let $A_{R_0}$ be the pole coefficient of
this pair, with the normalization preceding the common factor $1/(8sH)$.
Its $z$-cut finite part is
\begin{equation}
 F_{R_0+W,z}=E_{R_0+W}(s)
 -\frac{2s(-s^5-5s^4+18s^3+122s^2+151s-45)}{3(s^2+2s-3)}
 +A_{R_0}\ln\frac{s-1}{s}.
\end{equation}
The rational subtraction is the complete lower-basepoint constant.
The $W$ part has zero lower constant and zero soft pole. Finally, the
previously derived rational-sector matching shift converts this expression
to the common $(xy)^\eta$ convention. This order of operations preserves
all finite endpoint terms before combining the other real sectors.
